%%═══════════════════════════════════════════════════════════════════
%% Lecture 02 — Numerical Integration, Root Finding, and Linear Equations
%% 77315 — Computational Physics
%% Date: 12/04/2026
%%═══════════════════════════════════════════════════════════════════

\section{Lecture 02 — Numerical Integration, Root Finding, and Linear
  Equations}\label{sec:lec02}

\begin{center}
\begin{tabular}{ll}
  \textbf{Date:}\quad 12/04/2026 & \\
\end{tabular}
\end{center}
\hrule\vspace{1.5em}

%%──────────────────────────────────────────────────────────────────
\subsection{Numerical Integration (Quadrature)}
Numerical integration calculates the area under a curve $\int_a^b f(x)\,dx$
by evaluating the function at specific discrete points.

\begin{itemize}
    \item \textbf{Simpson's Rule}: Groups points in sets of three and fits a
    parabola through them. The area under this parabola gives:
    \begin{equation}
    \int_{-h}^{h} f(x)\,dx = \frac{h}{3}\bigl[f(-h) + 4f(0) + f(h)\bigr]
    \end{equation}
    It is remarkably accurate, possessing a global error of $\bigO{h^4}$.

    \item \textbf{Gauss-Legendre Integration}: Optimises both the weights
    $w_i$ and sampling locations $x_i$ simultaneously, doubling the degrees
    of freedom. The optimal $x_i$ match the roots of the Legendre polynomials.
    Using $n$ optimally placed points, Gauss-Legendre quadrature perfectly
    integrates any polynomial up to degree $2n-1$.

    \item \textbf{Recursive (Adaptive) Integration}: Evaluates the integral on
    a coarse interval, splits it in half, and re-evaluates. If the two answers
    agree within tolerance $\varepsilon$, it stops; otherwise it recursively
    dives deeper only in highly-varying sub-regions.
\end{itemize}

%%──────────────────────────────────────────────────────────────────
\subsection{Linear Equations (Gaussian Elimination)}
Solving a system of $n$ linear equations $Ax = b$ is the backbone of almost
all computational physics. \textbf{Gaussian Elimination} systematically
subtracts rows to transform $A$ into an upper-triangular matrix in
$\bigO{n^3}$ operations, then solves by \textbf{Back Substitution} in
$\bigO{n^2}$.

\begin{tcolorbox}[colback=green!5!white,colframe=green!50!black,
  title=The Critical Necessity of Partial Pivoting]
During Gaussian elimination, we divide the entire row by the pivot element
$A_{ii}$. If the pivot is very small, say $10^{-10}$, dividing amplifies
roundoff errors by a factor of $10^{10}$, destroying the rest of the matrix.

\textbf{Partial Pivoting}: Before processing each column, scan down the rows
to find the element with the largest absolute magnitude and swap it to the
top. By always dividing by the largest possible number, we suppress error
amplification and guarantee numerical stability.
\end{tcolorbox}

%%──────────────────────────────────────────────────────────────────
\subsection*{Recommended Reading}
\begin{itemize}
    \item \textit{Numerical Recipes}, W.H. Press et al.\ --- Ch.\ 4
    (Integration of Functions), Ch.\ 2.2 (Gaussian Elimination with
    Backsubstitution).
    \item \textit{Computational Physics}, Mark Newman --- Ch.\ 5 (Integrals
    and Derivatives), Ch.\ 6 (Solution of Linear and Nonlinear Equations).
\end{itemize}
